\documentclass[11pt]{article}
\usepackage[a4paper,margin=1in]{geometry}
\usepackage{amsmath,amssymb,amsthm,mathtools}
\usepackage{enumitem}

\newtheorem{theorem}{Theorem}[section]
\newtheorem{lemma}[theorem]{Lemma}
\newtheorem{proposition}[theorem]{Proposition}
\newtheorem{conjecture}[theorem]{Conjecture}
\newtheorem{remark}[theorem]{Remark}

\newcommand{\R}{\mathbb R}
\newcommand{\Om}{\Omega}
\newcommand{\Et}{E_{\hat t}}
\newcommand{\dT}{\delta_T}
\newcommand{\DI}{D_I}
\newcommand{\DHo}{D_H}
\newcommand{\A}{\mathcal A}
\newcommand{\HH}{\mathcal H}

\title{Attack on the Final Planar Radialization Step}
\author{}
\date{May 2026}

\begin{document}
\maketitle

\section{Setup and target}

Let \(E=\Et\subset\R^2\) be the \(\sqrt\delta\)-stopped level.  Normalize
\(|E|=\pi+O(\sqrt\delta)\), and write after harmless rescaling
\[
   |E|=\pi.
\]
The available inputs are
\[
   \dT(E)\le C\delta,
   \qquad
   \DI(E)\le \varepsilon,
   \qquad
   \DHo(E)\le \varepsilon,
   \qquad
   \varepsilon=C_K\sqrt\delta,
   \tag{1.1}\label{eq:inputs}
\]
and, from the selected-class interpolation,
\[
   \bigl\||\nabla u_E|-c_E\bigr\|_{L^\infty(\partial E)}
      \le \eta_K,
   \qquad
   \eta_K\to0 \quad(K\to\infty).
   \tag{1.2}\label{eq:serrin}
\]
We also use the uniform Lipschitz character of the stopped level, with
constant \(L\), and a bounded ambient diameter.

The needed conclusion is a small-amplitude radial graph \(G\) such that
\[
   |E\Delta G|\le C_K\sqrt\delta,
   \qquad
   \dT(G)\le C_K\delta .
   \tag{1.3}\label{eq:target}
\]
Once \eqref{eq:target} holds, the direct radial-graph theorem gives
\(\A(G)^2\le C\dT(G)\le C_K\delta\), and the transfer back to \(E\) and then
to \(U\) closes the desired estimate.

\section{Do not convexify everything}

The previous convex-envelope idea is useful but too crude if applied to all
of \(E\).  A wide shallow radial indentation can have area \(O(\sqrt\delta)\)
and perimeter excess \(O(\sqrt\delta)\), while already being a perfectly good
radial graph.  Convexifying it would create a volume error too large for the
torsion-deficit bookkeeping.

Therefore the repair should fill only the genuinely non-radial pockets:
holes, satellites, folds, same-ray missing intervals, and overhangs.  Radial
indentations with one crossing per ray should be left untouched.

\section{Components are ruled out by Serrin, not just by \(D_I\)}

The \(D_I\) estimate already gives satellite area \(O(\varepsilon^2)\).
In the selected Serrin situation one can say more.

\begin{lemma}[No small Serrin components]\label{lem:no-small-components}
Assume \(|E|=\pi\), \(P(E)\le2\pi+C\varepsilon\), and
\[
   c_0\le |\nabla u_E|\le c_1
   \quad\hbox{on }\partial E
\]
with \(c_0\) sufficiently close to \(1/2\).  Then, for \(\varepsilon\) small,
\(E\) has only one connected component.
\end{lemma}

\begin{proof}
On every connected component \(E_j\), the restriction of \(u_E\) is the
torsion function of \(E_j\).  The divergence theorem gives
\[
   |E_j|=\int_{\partial E_j}|\nabla u_E|\,d\mathcal H^1
      \ge c_0 P(E_j).
   \tag{3.1}
\]
The planar isoperimetric inequality gives
\[
   P(E_j)\ge 2\sqrt{\pi |E_j|}.
\]
If \(|E_j|>0\), then
\[
   |E_j|^{1/2}\ge 2c_0\sqrt\pi .
\]
For \(c_0\) close to \(1/2\), each nontrivial component has area close to
\(\pi\).  Since the total area is \(\pi\), there can be only one component
for \(\varepsilon\) and \(\eta_K\) small.
\end{proof}

\begin{remark}
This is a useful improvement over pure \(D_I\).  Tiny satellites are allowed
by small perimeter deficit, but they are incompatible with a pointwise Serrin
lower bound at the ball scale.
\end{remark}

\section{Closed holes are harmless}

Let \(p_h\) be the total perimeter of bounded hole boundaries.  Since holes
contribute to \(P(E)-2\pi\), \(p_h\le C\varepsilon\).  By the planar
isoperimetric inequality,
\[
   |H_{\rm holes}|
      \le \frac{p_h^2}{4\pi}
      \le C\varepsilon^2
      \le C_K\delta .
   \tag{4.1}\label{eq:holes-area}
\]
Filling them changes the volume by \(O(\delta)\).  The monotone
fill-and-rescale lemma from the earlier planar radialization note then gives
\[
   \dT(E_{\rm filled}^{\sharp})
      \le C(\dT(E)+|H_{\rm holes}|)
      \le C_K\delta .
   \tag{4.2}\label{eq:holes-fill}
\]

\section{The planar GMT pocket lemma}

Choose the eventual center \(z\) from the Brandolini/Plan-2 same-center
package.  The geometric input needed at this stage is not merely that
\(\partial E\) is Lipschitz.  One also needs the already advertised annular
core information, namely
\[
   B_{r_-}(z)\subset E\subset B_{r_+}(z),
   \qquad
   r_-=1-\eta,\quad r_+=1+\eta,
   \tag{5.1}\label{eq:core-annulus}
\]
with \(\eta\) smaller than the fixed graph-entry threshold.  Closed holes may
first be filled by Section~4; this changes area by \(O(\varepsilon^2)\), so
it is harmless for the present estimate.

For a.e. \(\theta\in S^1\), write
\[
   E_\theta=\{r>0:z+re^{i\theta}\in E\}.
\]
Since \(B_{r_-}(z)\subset E\), the slice \(E_\theta\) starts at the origin.
Define the radial hull
\[
   R(\theta):=\operatorname*{ess\,sup} E_\theta,
   \qquad
   G:=\{z+re^{i\theta}:0<r<R(\theta)\}.
   \tag{5.2}\label{eq:radial-hull}
\]
Thus one-crossing radial indentations are left unchanged.  The set
\(G\setminus E\) consists exactly of same-ray gaps lying behind some outer
piece of \(E\).

\begin{lemma}[Radial-hull boundary decomposition]\label{lem:hull-decomp}
Let \(E\subset\R^2\) be a bounded connected Lipschitz domain satisfying
\eqref{eq:core-annulus}.  Let \(G\) be defined by \eqref{eq:radial-hull}, and
let \(\{Q_j\}\) be the connected components of \(G\setminus E\).  Then \(G\)
has finite perimeter and
\[
   P(G)+\sum_j P(Q_j)\le P(E).
   \tag{5.3}\label{eq:decomposition}
\]
\end{lemma}

\begin{proof}
Translate \(z\) to the origin.  The polar map
\[
   \Phi:S^1\times(r_-,r_+)\to B_{r_+}\setminus \overline{B_{r_-}},
   \qquad
   \Phi(\theta,r)=re^{i\theta},
\]
is bi-Lipschitz, so rectifiability and perimeter identities may be checked in
the cylinder.  For a.e. \(\theta\), the one-dimensional slice \(E_\theta\) is
a finite-perimeter subset of \((0,\infty)\), and because of the core inclusion
its reduced boundary consists of an odd ordered family
\[
   r_1(\theta)<r_2(\theta)<\cdots<r_{2m(\theta)+1}(\theta).
\]
The outer endpoint \(r_{2m+1}(\theta)\) is \(R(\theta)\); the earlier
endpoints are precisely the two sides of the radial gaps that form
\(G\setminus E\).

We now apply the area formula, equivalently the BV graph-perimeter formula, to
the countably many polar sheets of \(\partial E\).  On a smooth sheet
\(\theta\mapsto r(\theta)\), the Euclidean length element is
\[
   ds=\sqrt{r^2\,d\theta^2+dr^2}.
\]
The same formula, with the usual vertical jump contribution, holds for BV
polar graphs.  The perimeter of \(G\) is obtained by selecting, for each
\(\theta\), the outermost endpoint \(R(\theta)\).  If this selected endpoint
jumps, the vertical jump segment in \(\partial G\) is exactly the jump term of
the corresponding outer sheet in the BV graph formula.

The components \(Q_j\) are the gaps between consecutive endpoint functions
\(r_{2i-1}(\theta)\) and \(r_{2i}(\theta)\).  Their perimeter is bounded by
the BV graph perimeters of these two endpoint sheets, including the vertical
jump terms at the angles where the gap opens or closes.  Consequently the
perimeter of \(G\), plus the perimeters of all pockets \(Q_j\), is bounded by
the sum of the polar graph perimeters of all sheets of \(\partial E\).  By the
area formula in the bi-Lipschitz polar chart, this sum is \(P(E)\), giving
\eqref{eq:decomposition}.
\end{proof}

\begin{theorem}[Lipschitz pocket lemma]\label{thm:pocket}
Let \(E\subset\R^2\) be a bounded connected Lipschitz domain with \(|E|=\pi\),
\[
   P(E)\le2\pi+\varepsilon,
\]
and assume the annular core condition \eqref{eq:core-annulus}.  Let \(G\) be
the radial hull \eqref{eq:radial-hull}.  Then \(G\) is a radial graph,
\[
   |G\setminus E|\le \frac{\varepsilon^2}{4\pi}.
   \tag{5.4}\label{eq:pocket-area}
\]
In particular, if \(\varepsilon=C_K\sqrt\delta\), then
\[
   |G\setminus E|\le C_K\delta .
   \tag{5.5}
\]
\end{theorem}

\begin{proof}
Write
\[
   v:=|G\setminus E|,\qquad h:=\sum_jP(Q_j).
\]
Since \(E\subset G\), \(|G|=\pi+v\).  By Lemma~\ref{lem:hull-decomp} and the
isoperimetric inequality applied to \(G\),
\[
   2\pi+\varepsilon
      \ge P(E)
      \ge P(G)+h
      \ge 2\sqrt{\pi(\pi+v)}+h
      \ge 2\pi+h .
   \tag{5.6}
\]
Thus \(h\le\varepsilon\).  Applying the planar isoperimetric inequality to
each pocket \(Q_j\) gives
\[
   |Q_j|\le \frac{P(Q_j)^2}{4\pi}.
\]
Summing over \(j\),
\[
   v
      =\sum_j |Q_j|
      \le \frac1{4\pi}\sum_jP(Q_j)^2
      \le \frac1{4\pi}\Bigl(\sum_jP(Q_j)\Bigr)^2
      \le \frac{\varepsilon^2}{4\pi}.
\]
This proves \eqref{eq:pocket-area}.  The inclusion
\eqref{eq:core-annulus} also gives
\[
   B_{r_-}(z)\subset G\subset B_{r_+}(z),
\]
so after the harmless area normalization \(G\) is a small-amplitude radial
graph whenever \(\eta\) is small.
\end{proof}

\begin{remark}
The theorem is deliberately about the radial hull, not the convex hull.  A
wide shallow one-crossing indentation has no same-ray gap and is not filled.
Thus the argument avoids the earlier \(O(\sqrt\delta)\)-volume loss.
\end{remark}

\section{Closure using the pocket lemma}

Fill the closed holes from Section~4 and then take the radial hull
\eqref{eq:radial-hull}.  Call the result \(G\).  The total filled volume is
\[
   |G\setminus E|\le C_K\delta .
   \tag{6.1}
\]
The monotone fill-and-rescale lemma yields
\[
   \dT(G^\sharp)
      \le C\bigl(\dT(E)+|G\setminus E|\bigr)
      \le C_K\delta .
   \tag{6.2}
\]
The annular amplitude required by the direct radial-graph theorem is exactly
\eqref{eq:core-annulus}.  Therefore \(G^\sharp\) satisfies
\eqref{eq:target}, and the direct radial-graph theorem gives
\[
   \A(G^\sharp)^2\le C\,\dT(G^\sharp)\le C_K\delta .
   \tag{6.3}
\]
Since \(|E\Delta G^\sharp|\le C_K\sqrt\delta\) after area normalization, the
usual transfer gives \(\A(E)\le C_K\sqrt\delta\), and then Step 5 transfers
the estimate back to \(U\).

\section{What remains outside this GMT step}

The GMT pocket step is proved above for Lipschitz boundaries.  The remaining
non-GMT inputs needed before applying it are:

\begin{enumerate}[label=\arabic*.]
\item the \(\sqrt\delta\)-stopped level;
\item exact deficit transfer \(\dT(E)\le C\delta\);
\item pointwise Serrin control on \(\partial E\);
\item no small components from the Serrin lower bound;
\item the annular core condition \eqref{eq:core-annulus} with the same center.
\end{enumerate}

Thus the old ``pocket area'' obstruction is no longer the bottleneck: once the
level is known to have the fixed annular core, radializing its same-ray gaps
costs \(O((P(E)-2\pi)^2)=O_K(\delta)\), which is the required torsion scale.

\end{document}
